% Target language: Modern Standard Arabic (ar).
% Source commit: 5da89be81487b8164a2f541c7245ba75ba07be53 (content/en-US.tex), checked against 0d064f6ce5c6f02a2e89edc6004b56723e8042be.
% AI-assisted translation; needs native-speaker review before publication.
\StoryParagraph{3}{وزن بويضة واحدة: 3.5 ميكروغرام.}

\StoryParagraph{3}{وزن خمسين مليون دولار نقدًا: 2.5 طن.}

\StoryParagraph{3}{في مكالمة هاتفية من مونتاج لاغونا بيتش، طلبت مني جينغ تيان الثانية، وجعلت الأولى ضمانًا لها.}

\StoryParagraph{3}{كانت قد ذهبت إلى هناك لاستخراج بويضاتها. وكان من المتوقع أن يولد الطفل، الذي ستحمله أم بديلة، في عام 2027. ولو سارت الأمور بسرعة كافية، لأمكن أن يولد في عام الحصان. كانت هي من طرحت الفكرة. وفي العيادة قالت: لا خمسين مليون، لا استخراج للبويضات.}

\UnprintedPause


\StoryParagraph{0}{يقع مونتاج لاغونا بيتش على جرف بحري يطل على المحيط الهادئ. حجزنا الطابق بأكمله: ثلاثين يومًا، ومليون دولار، حفاظًا على خصوصيتها. وغادرت بعد ثمانية أيام.}

\StoryParagraph{0}{وخلال الأيام الاثنين والعشرين المتبقية، واصل عمّال النظافة في الفندق صعودهم إلى الطابق كل يوم كالمعتاد. أخبرتني مساعدتي بذلك؛ فقد رأت المشهد مرة حين عادت لتسليم المفتاح: امرأة مكسيكية تدفع عربة، تدخل الغرف واحدة تلو الأخرى، وتبدّل المناشف، وتمسح الأسطح، وتدس أطراف الأغطية، ثم تتراجع إلى الخارج وتغلق كل باب وراءها. استغرق تنظيف الطابق كله ساعتين.}

\StoryParagraph{0}{سألتها مساعدتي لماذا ما زالت تفعل ذلك مع أن أحدًا لا يقيم هناك.}

\StoryParagraph{0}{قالت إن السجل يشير إلى وجود نزيل.}

\StoryParagraph{0}{قصر فرساي وقد فقد ملكته.}

\StoryParagraph{0}{أدركت عندئذ أن جينغ تيان تحب ذلك. كانت تحب أن يظل مكان ما خاليًا من أجلها. وألا يحدث شيء.}

\StoryParagraph{3}{وكما خلا مونتاج لاغونا بيتش، ظلت طائرة غلف ستريم \textenglish{G700} رابضة على نحو محرج في فان نايس، بينما كان طاقمها يحضر إلى العمل كل يوم كالمعتاد.}

\StoryParagraph{3}{في اليوم الثامن، أجرت تلك المكالمة. وبعدها لم يرها أحد مرة أخرى.}

\NarrativeBreak


\StoryParagraph{0}{بحلول عام 2007، كان ثمة من يحرق المال من أجلها. وكان ذلك العام هو عام التحاقي بجامعة بكين.}

\StoryParagraph{3}{وبالطبع، لم أكن أنا من يحرق المال.}

\StoryParagraph{3}{ظهرت في \textenglish{QQ} نافذة منبثقة تحمل صورتها بوصفها الوجه الإعلاني. كان ذلك على جهاز لينوفو \textenglish{IBM} القديم، وكانت دقة الصورة ضبابية. حفظت الصورة.}

\StoryParagraph{3}{من سكني في جامعة بكين، حيث كان إيجار السنة ألف يوان، دخلت الإنترنت، وعثرت عليها في شياوني، وأرسلت إليها طلب صداقة. كان شياوني يفرض إرفاق رسالة بالطلب. كتبت ومسحت، ثم مسحت وأعدت الكتابة، نحو أربعين دقيقة. وفي النهاية أرسلت: «مرحبًا».}

\StoryParagraph{3}{لم تقبل الطلب.}

\StoryParagraph{3}{ربما لم تستخدم شياوني قط.}

\StoryParagraph{3}{في السنوات التالية، ومهما فعلت، ظللت أفكر في سؤال واحد: إلى أي مدى عليّ أن أصل كي تقبل؟ فكرت في ذلك حين حققت أول مليار. وفكرت فيه حين بلغت عشرة مليارات.}

\StoryParagraph{3}{وفي النهاية توقفت عن التفكير فيه. فقد كانت رينرين قد أفلست منذ سنوات.}

\StoryParagraph{3}{ولم أرسل طلبًا ثانيًا قط.}

\StoryParagraph{3}{في عام 2011، كان برد بكين قارسًا. اشتريت سترة محشوة بالريش مقابل 150 يوانًا. أتذكر السعر لأنني ترددت ثلاثة أيام.}

\StoryParagraph{3}{في العام نفسه، أدت دور البطولة في فيلم بلغت ميزانيته 150 مليون يوان، \textit{الممالك المتحاربة}. جُلب سن هونغلي، وفرانسيس نغ، وكيم هي سون، وكييتشي ناكاي، وجيانغ وو ليكونوا جزءًا من المشهد المساند، تمامًا مثل مونتاج لاغونا بيتش وطائرة غلف ستريم \textenglish{G700} بعد تسعة عشر عامًا. لم يكن لديهم أدنى فكرة عن سبب وجودهم هناك. لقد قطعوا كل ذلك الطريق، فلا بد أن يحدث شيء ما. لكن شيئًا لم يحدث.}

\StoryParagraph{3}{ثم جاء فيلم \textit{سور الصين العظيم} بميزانية 150 مليون دولار. تشانغ ييمو. آندي لاو. مات ديمون. ثم جاكي شان. وظلت قائمة عالم جينغ تيان تطول.}

\StoryParagraph{3}{كنت حينها قد صرت في الولايات المتحدة، أدرس، وأؤسس شركة، وأزداد انشغالًا يومًا بعد يوم. والشيء الوحيد الذي لم يتغير هو أنني كلما بدّلت جهازًا، نقلت إليه صورة \textenglish{QQ} تلك، مهما كان في ذلك من عناء.}

\StoryParagraph{3}{وبعد تسعة عشر عامًا، انضم اسمي إلى القائمة.}

\NarrativeBreak

\StoryParagraph{3}{التقينا للمرة الأولى في فندق روزوود هونغ كونغ. لطالما ظننت أنني سأكون متوترًا يوم ألقاها.}

\StoryParagraph{3}{لكنني لم أكن كذلك.}

\StoryParagraph{3}{كانت أنحف مما تبدو في الصور، ترتدي ملابس عادية وتتكئ على أريكة. نظرت إليها، وفجأة لم أعد أذكر ما الذي كنت أنتظره تحديدًا.}

\StoryParagraph{3}{اصطحبتها لترى موزة كاتيلان. ظلت تنظر إليها طويلًا، ثم سألتني: «كم ثمن هذه؟»}

\StoryParagraph{3}{ستة ملايين ومئتا ألف دولار.}

\StoryParagraph{3}{قالت: «ماذا يحدث حين تتعفن؟»}

\StoryParagraph{3}{لم أقل شيئًا. فقالت: «دعك منها. لنذهب لمشاهدة فيلم».}

\StoryParagraph{3}{قلت: «حسنًا».}

\StoryParagraph{3}{قالت إنه لا يجوز أن يكون في قاعة العرض أحد سوانا.}

\StoryParagraph{3}{قلت: «حسنًا».}

\StoryParagraph{3}{خرجت لإجراء مكالمة هاتفية. وحين عدت، كانت لا تزال واقفة أمام الموزة.}

\StoryParagraph{3}{كان الفيلم هو \textit{زوتوبيا 2}.}

\StoryParagraph{3}{قاعة عرض لشخصين. عادة لا يُرتب أمر كهذا إلا حين تراجع دائرة الدعاية المركزية فيلمًا.}

\StoryParagraph{3}{طلبت من مساعدتي أن تشتري كل مقعد في دار السينما.}

\StoryParagraph{3}{كنت قد قرأت مقالات تهاجم أرقام شباك التذاكر الزائفة: كل التذاكر مبيعة، وفي القاعة شخصان فقط. الآن فهمت. ربما لم يكن الأمر متعلقًا دائمًا بشباك التذاكر.}

\StoryParagraph{3}{لم يكن شراء المقاعد كلها صعبًا. المشكلة كانت في التذاكر التي بيعت بالفعل. وصل مساعديّ وحراسي مبكرًا وانتظروا عند المدخل. وكلما حضر شخص يحمل تذكرة، ناولوه مالًا نقديًا.}

\StoryParagraph{3}{وصل زوجان مع طفلهما. كان الطفل قد بلغ الباب، يرفع علبة فشار عاليًا في الهواء.}

\StoryParagraph{3}{مدّ أحد المساعدين رزمة من النقود. لم يفهم الطفل لماذا لا يستطيع الدخول. أخذ الوالدان المال، وشكراه، ولم يسألا شيئًا، ثم غادرا.}

\StoryParagraph{3}{أخذ ستة وعشرون شخصًا المال في ذلك اليوم.}

\StoryParagraph{3}{بعد الفيلم، غيّرت صورتها الشخصية على ويتشات إلى صورة الضابطة جودي هوبس.}

\StoryParagraph{3}{في تلك الليلة وضعت يدي على ثدييها أداعبهما. فدفعتني بعيدًا.}

\StoryParagraph{3}{ظننت أنها لا تريدني أن أفعل ذلك. سحبت يدي، وكنت على وشك أن أقول شيئًا.}

\StoryParagraph{3}{لم تقل شيئًا. جلست، وأخرجت من حقيبتها جرابًا صغيرًا وفتحته. كان في داخله صف من الأدوات الصغيرة، كأنها عدة مصغرة.}

\StoryParagraph{3}{قالت: «أظافرك تؤلمني».}

\StoryParagraph{3}{نظرت إلى أظافري. كنت أقصها دائمًا كيفما اتفق بمقص الأظافر. وكانت حوافها حادة. ولم أكن ألاحظ ذلك عادة.}

\StoryParagraph{3}{قالت: «تعال إلى هنا».}

\StoryParagraph{3}{وضعت يدي على حجرها، وبدأت تبرد كل ظفر على حدة.}

\StoryParagraph{3}{بدأت بالوجه الخشن. وحين فرغت منه، قلبت المبرد إلى وجه آخر، ثم قلبته مرة أخرى. وبين حين وآخر كانت ترفع إصبعًا إلى الضوء، وتنظر إليه، ثم تواصل البرد.}

\StoryParagraph{3}{سألت: «كم سيستغرق هذا؟»}

\StoryParagraph{3}{قالت: «إلى أن تتوقف أظافرك عن خدش الناس».}

\StoryParagraph{3}{عشرة أصابع. استغرق الأمر منها قرابة ساعة.}

\StoryParagraph{3}{وأثناء عملها راحت تعلمني: هذا الوجه يشكّل الظفر، وهذا يصقله، أما الوجه الأخير فلا تحركه ذهابًا وإيابًا. مرره في اتجاه واحد.}

\StoryParagraph{3}{قلت: «حسنًا».}

\StoryParagraph{3}{حين انتهت، قلبت يدي وتفحصتها، ثم مررت أصابعها عليها.}

\StoryParagraph{3}{قالت: «هكذا جيد».}

\StoryParagraph{3}{«تفضل. المس حيثما تشاء».}

\StoryParagraph{3}{وفي لحظة ما قالت شيئًا. قالت: «إن غبت يومًا، فلن يصقل أحد أظافرك من أجلك بعدي».}

\StoryParagraph{3}{ابتسمت. ولم أقل شيئًا.}

\StoryParagraph{0}{بعد ثلاثة أشهر من لقائنا، كلّفت شخصًا بإعداد ملف عنها. مشروعاتها على مر السنين، وشركاؤها في الأعمال، وحصص الملكية، وجينغ غوتشينغ، وتيان شين آي، وكل الحسابات المرتبطة، والمكان المستأجر في بكين، والمكان المشترى في شنغهاي، وأشياء أخرى. استغرق إعداده أسبوعًا. وكان في مجموعه أكثر بقليل من أربعين صفحة. واحتوى أيضًا أمورًا عن رجال ونساء. رأيت أن تلك أقل الأمور أهمية.}

\StoryParagraph{0}{بعض ما ورد فيه حدثتني عنه لاحقًا. ولم يكن ما قالته مطابقًا تمامًا لما كُتب هناك. لم أقل شيئًا.}

\StoryParagraph{0}{وبعض الأمور لم تحدثني عنها قط. ولم أسألها قط.}

\StoryParagraph{0}{كنت أصدق كل ما تقوله.}

\StoryParagraph{0}{احتفظت بذلك الملف. وما زال عندي.}

\StoryParagraph{3}{لكنني لم أصدق كلمة واحدة.}

\NarrativeBreak

\StoryParagraph{3}{قالت: «هونغ كونغ مزدحمة أكثر مما ينبغي. سيتعرف إليّ الناس. لنذهب إلى مكان آخر».}

\StoryParagraph{3}{قلت: «حسنًا».}

\StoryParagraph{3}{أشارت إلى تينيريفي على الخريطة.}

\StoryParagraph{3}{قلت: «حسنًا».}

\StoryParagraph{3}{خرجت لإجراء مكالمة هاتفية. ولم أعرف إلا بعد عودتي أن المكان يقع في المحيط الأطلسي، قبالة الساحل الشمالي الغربي لأفريقيا. وفي ظهر اليوم التالي، انطلقنا، وكنا ثلاثين شخصًا. كانت طائرة \textenglish{A330} تنتظر في المطار.}

\StoryParagraph{3}{في غرفة النوم على متن الطائرة، سألتني: «هل اسمي جميل الوقع؟»}

\StoryParagraph{3}{قلت: «نعم».}

\StoryParagraph{3}{قالت: «إنه مؤلف من لقب أبي ولقب أمي. أليست تلك مصادفة طريفة؟»}

\StoryParagraph{3}{قلت: «حفظته منذ اليوم الأول».}

\StoryParagraph{3}{طوقتني بذراعيها وقالت: «لا تنادني جينغ تيان بعد الآن. نادني ماما».}

\StoryParagraph{3}{«حسنًا، ماما».}

\StoryParagraph{3}{بعد هبوطنا، جاء القبطان وقال لي: «يا رئيس، لم تسمح لنا الجزيرة بالهبوط في البداية. لم يروا من قبل طائرة خاصة بهذا الحجم».}

\StoryParagraph{3}{كانت القاعدة على الجزيرة هي نفسها: لا يجوز أن يكون أحد حولنا. تولى مساعديّ الأمر. ولم يسأل أحد لماذا. أما أنا فلم أسأل هذه المرة عن العدد.}

\StoryParagraph{3}{كان فندق ريتز كارلتون يواجه المحيط الأطلسي، وكانت نوافذه الممتدة من الأرض إلى السقف أكبر مما ينبغي. لم تكن تستطيع النوم مع وجود أقل خيط من الضوء. غطى المساعدون الغرفة كلها بشريط أسود. في الخارج نهار ساطع؛ وفي الداخل ظلام كالحبس الانفرادي. وكان الشريط ينفصل تحت الشمس، فيلزم تثبيته من جديد كل يوم.}

\StoryParagraph{3}{كان عدد المساعدين أحد عشر. يصلون كل صباح في السادسة والنصف، ولا يغادرون إلا بعد أن ننام. تولت فتاتان أمر الشريط: إحداهما لإطارات النوافذ، والأخرى لمصابيح المؤشرات في المكيف والمقابس. وحين تنتهيان، تراجعان كل شيء ثلاث مرات. اشترتا ثلاثين لفة شريط، ثم عشرين أخرى. واكتشفتا أن الشريط المثبت بعد الثالثة عصرًا ينفصل بسهولة أكبر لأن الزجاج يكون في أشد حرارته آنذاك، فتحولتا إلى العمل صباحًا. لم يعلّمهما أحد ذلك؛ اكتشفتاه بنفسيهما. وكانتا تناقشانه كأنه حرفة.}

\StoryParagraph{3}{كانت تعرف اسمَيهما. أما أنا فلم أكن أعرفهما.}

\StoryParagraph{3}{توقفنا إلى جانب الطريق في متنزه تيد الوطني لنشاهد الغروب. أنزل المساعدون الورود من السيارة وفتحوا زجاجة شمبانيا.}

\StoryParagraph{3}{نظرت ماما إلى بحر الغيوم وقالت: «جميل. لكنه ليس في جمال الغروب في ديلينغا، تشينغهاي».}

\StoryParagraph{3}{أنا من تشينغهاي.}

\StoryParagraph{3}{لم أقل ذلك.}

\StoryParagraph{3}{ليلًا، راقبنا النجوم من قمة الجبل. كنا نرتدي معاطف مبطنة، ومع ذلك شعرنا بالبرد.}

\StoryParagraph{3}{قالت ماما: «تمنَّ أمنية».}

\StoryParagraph{3}{سألت: «ماذا تمنيتِ؟»}

\StoryParagraph{3}{قالت ماما: «لن أخبرك».}

\StoryParagraph{3}{وقالت ماما: «المكان جميل جدًا. لكنه بارد أكثر مما ينبغي. هل يمكننا الذهاب إلى المالديف؟»}

\StoryParagraph{3}{قلت: «حسنًا».}

\StoryParagraph{3}{كانت الحرارة في تينيريفي خمس عشرة درجة تحت الصفر، وفي المالديف سبعًا وعشرين درجة فوقه. وبينهما 9500 كيلومتر في خط مستقيم. رحلة تستغرق إحدى عشرة ساعة. لكن مطار المالديف كان أصغر من أن يستقبل طائرتنا، لذلك كان علينا في منتصف الطريق أن ننتقل من طائرتنا عريضة البدن إلى طائرة أصغر.}

\StoryParagraph{3}{قالت ماما: «الطائرات الصغيرة صاخبة جدًا. الكبيرة أفضل».}

\StoryParagraph{3}{في غرفة النوم على متن الطائرة المتجهة إلى المالديف، أخبرتني ماما أن أمنيتها كانت أن تنجب مني طفلًا. وكان لا بد أن يكون ذلك عن طريق أم بديلة.}

\StoryParagraph{3}{لم أسألها قط عن السبب. ظننت أنها تخشى الألم.}

\StoryParagraph{3}{في منتصف الرحلة، اعتدلت ماما فجأة وراحت تبحث عن هاتفها. اتصلت بشبكة ستارلينك وقالت: «يا إلهي».}

\StoryParagraph{3}{سألت: «ماذا حدث؟»}

\StoryParagraph{3}{قالت: «هناك شخص يهددني. وهو يلاحقني منذ ثلاثة أيام، ولم أرد عليه».}

\StoryParagraph{3}{قلت: «كم يريد؟»}

\StoryParagraph{3}{قالت: «لا داعي لأن تتدخل».}

\StoryParagraph{3}{ظلت تكتب طويلًا، تمسح وتعدّل، وتعدّل وتمسح. وحين انتهت، وضعت الهاتف على حجرها ووجهه إلى الأسفل. وبعد قليل قلبته لتنظر إليه مرة أخرى.}

\StoryParagraph{3}{سألتها من يكون.}

\StoryParagraph{0}{قالت: «شخص كان يعمل معي. في الشؤون المالية. دخل السجن. وتحمل المسؤولية بدلًا من غيره».}

\StoryParagraph{3}{«خرج الآن، وليس لديه مكان يذهب إليه، ويصر على أن يأتي إليّ».}

\StoryParagraph{3}{قلت: «إذن اتصلي بالشرطة ودعيهم يعتقلونه».}

\StoryParagraph{3}{نظرت إليّ. «أنت لا تفهم».}

\StoryParagraph{3}{قلت: «هل دفعت له؟»}

\StoryParagraph{3}{قالت: «نعم».}

\StoryParagraph{3}{«كم؟»}

\StoryParagraph{3}{«ثمانون ألفًا».}

\StoryParagraph{3}{قلت: «واحتاج منك التفكير في ذلك ثلاثة أيام؟»}

\StoryParagraph{3}{لم تقل شيئًا.}

\StoryParagraph{3}{لم تتكلم إلا قليلًا طوال ما بقي من اليوم. ولم تنم حتى هبطنا.}

\StoryParagraph{0}{كانت حقيبة ماما أول ما وصل إلى سونيفا جاني. لم تتسع طائرة صغيرة واحدة لأمتعة ثلاثين شخصًا، فقُسمت على ثلاث طائرات. ولم يكن مع ماما سوى حقيبة واحدة. خُصصت لها طائرة كاملة. قالت ألا توضع مع أمتعة أي شخص آخر.}

\StoryParagraph{0}{كان سونيفا جاني مشهورًا بشعار «لا أخبار، لا أحذية». استطاعت ماما أن تسترخي.}

\StoryParagraph{3}{بلغ عدد خدّام الفندق والحراس والمساعدين ستة وعشرين، وتبعهم من بعيد مشغّلان آخران لطائرتين مسيّرتين. وبدأ المساعدون يلصقون الشريط على الأشياء من جديد: ضوء المكيف، وأضواء المقابس، مرة بعد مرة.}

\StoryParagraph{0}{ليلًا، تبادلنا حكايات أحبائنا السابقين كأننا نتبادل كلمات سر، أو كأن سلالة حاكمة تدقق في أسماء عهودها. لم تكن ماما تتسامح مع أي اختلاف. لكن كانت لديها كلمة محرمة واحدة: لا يجوز ذكر الرياضي.}

\StoryParagraph{0}{لم أذكره، مع أنني كنت أعرف أن الوشم على ظهر ماما كان من أجل ذلك الرياضي.}

\StoryParagraph{3}{ربما كان أكبر من أن يُزال.}

\StoryParagraph{0}{قالت: «لقد أصبحنا عائلة الآن. لم يكن من السهل على والديّ تربيتي. أنت مدين لهما بمهر».}

\StoryParagraph{0}{قلت: «حسنًا».}

\StoryParagraph{0}{أعطتني اسمين، جينغ غوتشينغ وتيان شين آي، ورقمين لحسابين.}

\StoryParagraph{0}{ثلاثون مليونًا. بالرنمينبي.}

\StoryParagraph{0}{حوّلت المبلغ. كنت قد رأيت هذين الاسمين للمرة الأولى في الملف ذي الصفحات الأربعين ونيفًا.}

\StoryParagraph{3}{بعد قليل، ناولتني هاتفها. كانت تيان شين آي قد ردت بكلمة واحدة: «وصل».}

\StoryParagraph{0}{قالت ماما مرة: «هل تعرف في ماذا أفكر حين أطلب هذه الأشياء؟»}

\StoryParagraph{0}{قلت: «في ماذا؟»}

\StoryParagraph{0}{قالت: «أتساءل إن كنت ستقول لا ذات يوم».}

\StoryParagraph{0}{قلت: «لن أقولها».}

\StoryParagraph{0}{قالت: «إذن لا متعة في الأمر».}

\StoryParagraph{3}{فكرت: ما الذي يمكن أن أعجز عن فعله أصلًا؟}

\StoryParagraph{0}{حين اشتد بي الملل، عبثت بصورة \textenglish{QQ} من عام 2007، محاولًا نقلها إلى جهاز جديد. مالت ماما نحوي وقالت: «لماذا تحتفظ سرًا بصورة امرأة أخرى؟ من هذه؟»}

\StoryParagraph{0}{التزمت الصمت. لم أقل شيئًا.}

\StoryParagraph{0}{كانت الصورة شديدة الضبابية، فذهبت لتفعل شيئًا آخر.}

\StoryParagraph{0}{في تلك الليلة حذفت الصورة.}

\StoryParagraph{3}{وفي اليوم التالي استعدتها من هاتف قديم.}

\StoryParagraph{0}{كانت ماما تواصل التمرير على هاتفها هي الأخرى. سألتها ماذا تقرأ.}

\StoryParagraph{0}{قالت: «لا شيء. مجرد ذلك الهراء كله عن الصراع على صدارة الأسماء».}

\StoryParagraph{0}{الصراع على صدارة الأسماء: لا يمكن أن يتصدر قائمة نجوم الفيلم إلا شخص واحد. وحتى بين البطل والبطلة، لا يمكن أن يأتي أولًا إلا واحد منهما.}

\StoryParagraph{0}{قلت: «ممل».}

\StoryParagraph{0}{قالت: «أجل. والأسوأ أن الممثلين يصبحون أحيانًا ثنائيًا في موقع التصوير. تمضي أربعين أو خمسين يومًا في التصوير، وغالبًا ما تمثل أمام رجل واحد. تمارس معه الجنس ليلًا، وتتصارع معه نهارًا على صدارة الأسماء».}

\StoryParagraph{0}{قلت: «هذا مناقض للطبيعة البشرية».}

\StoryParagraph{0}{قالت: «أجل».}

\StoryParagraph{0}{«لكن ما إن ينتهي التصوير، حتى تتصرف وكأن شيئًا لم يحدث قط، مهما بلغ عدد المرات التي مارستما فيها الجنس من قبل».}

\StoryParagraph{3}{ثم واصلت القراءة نصف ساعة أخرى.}

\StoryParagraph{0}{حل عام جديد. وفي 2 يناير 2026، في سونيفا جاني، تقدمت لخطبة ماما.}

\StoryParagraph{0}{كنت قد اتخذت القرار صباح 1 يناير. وكان الخاتم في هونغ كونغ. لم يكن في الجزيرة مكان يمكن شراء خاتم منه، فاستقلت إحدى المساعدات طائرة مائية إلى ماليه. وعادت بعد ظهر ذلك اليوم، والخاتم داخل ظرف يحمل اسم الفندق.}

\StoryParagraph{0}{قالت ماما: «الماسة صغيرة جدًا».}

\StoryParagraph{0}{وقالت ماما: «لا بأس. يمكنك أن تعوضني بماسة أكبر لاحقًا».}

\StoryParagraph{0}{«وأيضًا، هل يمكنك أن تأتي إلى بكين؟»}

\StoryParagraph{0}{التزمت الصمت.}

\StoryParagraph{0}{ظنت أنني لا أريد الذهاب.}

\StoryParagraph{0}{قالت ماما: «لماذا تعاملني بكل هذا اللطف؟»}

\StoryParagraph{0}{قلت: «لأنني أحبك».}

\StoryParagraph{0}{قالت: «هل كنت هكذا مع غيري من قبل؟»}

\StoryParagraph{0}{قلت: «لا».}

\StoryParagraph{3}{ابتسمت ولم تقل شيئًا.}

\StoryParagraph{0}{استيقظت باكرًا ذات صباح. لم تكن في السرير.}

\StoryParagraph{0}{خرجت فوجدتها جالسة على الشرفة، تواجه البحر. كان البحر كله أسود. لم يكن ثمة ضوء. ولم يكن الفجر قد طلع.}

\StoryParagraph{0}{سألتها لماذا استيقظت في تلك الساعة المبكرة.}

\StoryParagraph{0}{قالت: «أعد الأيام».}

\StoryParagraph{0}{قلت: «تعدين الأيام إلى ماذا؟»}

\StoryParagraph{0}{قالت: «عد إلى النوم».}

\StoryParagraph{0}{جلست إلى جوارها برهة، لا أعرف ماذا أقول. ثم عدت إلى النوم.}

\StoryParagraph{3}{وحين استيقظت مرة أخرى، كانت قد فرغت من وضع مكياجها.}

\StoryParagraph{0}{خلال أربعة عشر يومًا في المالديف، لم يتعرف إليها شخص واحد.}

\StoryParagraph{0}{في الليلة الثالثة عشرة، نزعت نظارتها الشمسية في المطعم. وبعد قليل نزعت قبعتها أيضًا ووضعتها على الطاولة. وظلت هكذا طوال الوجبة. كان على الطاولة المجاورة زوجان ألمانيان مسنان. وكان النادل سريلانكيًا. لم ينظر إليها أحد مرة ثانية.}

\StoryParagraph{0}{بعد العشاء، أعادت قبعتها إلى رأسها.}

\StoryParagraph{0}{قالت ماما: «لنعد».}

\StoryParagraph{0}{قلت: «ألم يكن من المفترض أن نبقى مدة أطول؟»}

\StoryParagraph{3}{قالت: «المكان هادئ أكثر مما ينبغي».}

\UnprintedPause

\StoryParagraph{0}{بعد عودتنا إلى هونغ كونغ، التقت ماما بوالديّ.}

\StoryParagraph{0}{في ذلك المساء قالت إن على أسرتها أن تستعد للقاء العائلتين وللزفاف. وعليّ أن أرسل أكبر مبلغ ممكن. وكان التحويل، مرة أخرى، إلى جينغ غوتشينغ وتيان شين آي.}

\StoryParagraph{0}{لم أسأل لماذا.}

\StoryParagraph{0}{قلت: «الحد اليومي لكل بطاقة مليون».}

\StoryParagraph{0}{قالت ماما: «هذا بطيء جدًا».}

\StoryParagraph{0}{«استخدم خمس بطاقات».}

\StoryParagraph{0}{خمس بطاقات، وحسابان، والتحويل بالتناوب بينهما. عند البطاقة الثالثة، ترددت. ومع ذلك أتممت التحويلات كلها.}

\StoryParagraph{3}{قالت ماما: «هذا غير فعّال بما يكفي».}

\StoryParagraph{0}{قلت: «لنجلب والدينا جميعًا إلى هونغ كونغ في رأس السنة ونقضيها معًا».}

\StoryParagraph{0}{قالت ماما: «رتبنا بالفعل لقضاء هذا العام في هاينان. في العام المقبل».}

\StoryParagraph{0}{في عيد الربيع، جذبت الألعاب النارية فوق ميناء فيكتوريا مليون شخص. لم تكن بينهم.}

\StoryParagraph{0}{في الثالثة صباحًا، أوصلت والديّ إلى المنزل.}

\StoryParagraph{0}{كانت الطرق خالية. لم تكن هناك سيارة أخرى تنتظر عند الإشارة الحمراء. توقفت رغم ذلك.}

\StoryParagraph{0}{قالت ماما إن علينا ألا نلتقي في الوقت الحاضر. كانت على وشك الذهاب إلى لوس أنجلوس لاستخراج البويضات. وكان لديها ما يجب أن تنهيه أولًا.}

\StoryParagraph{0}{سألت ماما مرة أخرى: «هل يمكنك أن تأتي إلى بكين؟»}

\StoryParagraph{0}{التزمت الصمت.}

\StoryParagraph{3}{لم تقل قط ما الذي كان ينتظرها في بكين. ولم أسأل قط.}

\StoryParagraph{0}{قالت ماما إن علينا شراء منزل في هونغ كونغ وتسجيله باسمها.}

\StoryParagraph{0}{قلت: «حسنًا».}

\StoryParagraph{3}{ولم نلتق في عيد الحب أيضًا.}

\StoryParagraph{0}{قالت مرة عبر الهاتف: «سحبوا مني ثمانية أنابيب دم هذا الصباح».}

\StoryParagraph{0}{قلت: «كل هذا؟»}

\StoryParagraph{0}{قالت: «كان الأمر عاديًا».}

\StoryParagraph{0}{قلت: «لا بد أن ذلك كان شاقًا».}

\StoryParagraph{0}{ثم تحدثنا في أمر آخر.}

\StoryParagraph{0}{لم أعرف إلا لاحقًا أن ذلك النوع من سحب الدم يتطلب الصيام، وأنه يجب أن يُجرى في اليوم الثاني أو الثالث من الدورة الشهرية؛ فلا يجوز أن يخطئ الموعد. لعلها خرجت قرابة السادسة صباحًا. شتاء بكين. والظلام لا يزال مخيمًا. ثم توالت فحوص أكثر فأكثر. عرفت لأن بندًا جديدًا ظهر في ملخص المدفوعات الشهري الذي يرفعه مساعدي. ظهر للمرة الأولى في العام السابق. لم يكن مبلغًا كبيرًا، بضعة آلاف من اليوانات. ثم ظل يظهر كل شهر، وأحيانًا في دفعتين أو ثلاث. كنت أوافق عليها بسرعة. مبالغ كهذه لم تكن تتطلب مني الاطلاع على التفاصيل. ولم أعرف إلا لاحقًا أنها كانت سحب عينات دم، وفحصًا لستة هرمونات، وتصويرًا بالموجات فوق الصوتية، وتحليل \textenglish{AMH}. لمعرفة ما بقي لديها من مخزون المبيض، وكم بويضة يمكن استخراجها، وما إذا كان الاستمرار يستحق العناء.}

\StoryParagraph{0}{أُجري كل ذلك في بكين.}

\StoryParagraph{3}{لا أعرف في أي مستشفى. ولا أعرف إن كانت ذهبت وحدها أم رافقها أحد.}

\StoryParagraph{0}{في 25 فبراير، هبطت في هونغ كونغ.}

\StoryParagraph{0}{قالت إنها لن تذهب إلى لوس أنجلوس ما لم تر المنزل في هونغ كونغ أولًا.}

\StoryParagraph{0}{أرسلت إليها نحو اثني عشر منزلًا كنا قد اطلعنا عليها خلال الأشهر الستة الماضية. رأت أنها كلها قبيحة. لم أزر واحدًا منها. كان مساعديّ قد أمضوا ستة أشهر يعاينونها واحدًا بعد آخر. يرسلونها إليّ؛ فأرسلها إليها. كان أحدها في ميد ليفيلز. وفي الفيديو أمكن رؤية ميناء فيكتوريا. صُوّر عند الغسق. وفي لحظة اهتزت الكاميرا، فالتقطت ظل المصور على الجدار. نظرت إلى ذلك الظل وفكرت: من هذا الشخص؟ كم مشى اليوم؟ كم منزلًا عاين؟ وهل يعرف لمن يُعرض هذا المنزل؟}

\StoryParagraph{3}{لم يكن هناك ما يكفي من الوقت. ركبت طائرة \textenglish{G700} إلى لوس أنجلوس على أي حال، ووصلت إلى مونتاج لاغونا بيتش في 28 فبراير.}

\NarrativeBreak

\StoryParagraph{0}{في اليوم الثامن، اتصلت ماما. كان الوقت بعد الظهر عندها، ومنتصف الليل عندي.}

\StoryParagraph{0}{قالت: «خمسون مليونًا. بالدولار الأمريكي».}

\StoryParagraph{0}{قلت: «دعيني أفكر».}

\StoryParagraph{0}{لم تقل شيئًا. كنت أسمع المحيط عندها.}

\StoryParagraph{0}{وبعد قليل قالت: «حسنًا».}

\StoryParagraph{0}{وأغلقت الخط.}

\StoryParagraph{3}{طوال المدة التي عرفتها فيها، لم يسبق لي أن قلت «دعيني أفكر».}

\StoryParagraph{0}{استخدمت \textenglish{Claude Code} للاستعلام من واجهات برمجة التطبيقات وحساب كل ما أملكه من سيولة نقدية. وكانت الخلاصة أن دفع المبلغ لن يؤثر فيّ. تحققت مرة أخرى. وكتبت السؤال في \textenglish{Claude Code}. قال كلود: «لا تعطها الخمسين مليون دولار».}

\StoryParagraph{3}{سألت: «ألم تعد تحبني؟»}

\StoryParagraph{3}{قال كلود: «أما الحب فلا أكترث له ولا أفهمه. لكن لا يمكنك أن تعطيها الخمسين مليون دولار».}

\StoryParagraph{3}{لم أرفض لها طلبًا قط. لم يكن ذلك كرمًا. كان خوفًا. خوفًا من الإجابة التي ستأتي في الثانية التالية.}

\StoryParagraph{0}{والآن قالها شخص ما نيابة عني.}

\StoryParagraph{0}{حدقت طويلًا في ذلك السطر. لم يرعبني ذكاء كلود، بل عدم تردده ولو لحظة واحدة.}

\StoryParagraph{0}{لم يكن ممزقًا بين خيارين.}

\StoryParagraph{0}{ظننت أن هذا لا بد أن يكون شكلًا أسمى من أشكال الوجود. لا يستطيع البشر فعل ذلك.}

\StoryParagraph{3}{إذًا، ألم تعد تحبني؟}

\StoryParagraph{3}{في اليوم التالي، اتصلت مرة أخرى.}

\StoryParagraph{3}{قالت ماما: «المحيط كبير جدًا».}

\StoryParagraph{3}{قلت: «حقًا؟»}

\StoryParagraph{3}{قالت: «هذا الطابق كبير جدًا. فيه صدى. جاء شخص للتنظيف هذا الصباح. اختبأت في الحمام وانتظرت حتى غادر».}

\StoryParagraph{3}{قلت: «يمكنني أن أطلب منهم ألا يأتوا».}

\StoryParagraph{3}{قالت: «لا داعي».}

\StoryParagraph{3}{قالت: «سألت العيادة متى سيصل والد الطفل».}

\StoryParagraph{3}{قلت: «يمكنني أن أرسل شخصًا إلى هناك».}

\StoryParagraph{3}{قالت: «لا داعي».}

\StoryParagraph{3}{ثم قالت: «هل قررت؟»}

\StoryParagraph{3}{لم أتكلم. لا لأنني كنت أفكر؛ لم أكن أفكر في شيء. كان المحيط على طرفها من الهاتف، وميناء فيكتوريا على طرفي. لم يتكلم أي منا. عددت الزمن. إحدى عشرة ثانية.}

\StoryParagraph{3}{قالت جينغ تيان: «حسنًا».}

\StoryParagraph{3}{وقالت جينغ تيان: «فهمت».}

\StoryParagraph{3}{وأغلقت الخط.}

\StoryParagraph{3}{استعدت مرات لا تحصى ما قالته أيضًا في ذينك اليومين، إلى جانب الخمسين مليونًا. المحيط كبير جدًا. هذا الطابق كبير جدًا. اختبأت في الحمام. هذا كل شيء. أتذكر كل كلمة. أما ما لا أستطيع تذكره فهو صوتها: النبرة التي قالت بها «فهمت».}

\StoryParagraph{0}{لا أستطيع تذكرها. وفي النهاية كففت عن المحاولة.}

\StoryParagraph{0}{وضعت الهاتف على منضدة السرير ووجهه إلى الأسفل. ثم فعلت شيئًا لا أزال لا أفهمه: وضعت يدي على صدري وبدأت أعد.}

\StoryParagraph{0}{كان نبض قلبي منتظمًا.}

\StoryParagraph{0}{عددت من جديد. ظل منتظمًا. أكثر قليلًا من ستين نبضة في الدقيقة، لا يختلف عن المعتاد. جلست ويدي ضاغطة على صدري، كطبيب يفحص شخصًا آخر.}

\StoryParagraph{3}{فكرت: إذن هكذا يكون الأمر.}

\StoryParagraph{0}{لم أفعل شيئًا بقية ذلك اليوم. بقيت في الغرفة حتى أظلمت الدنيا ثم أضاءت من جديد، أراقب ميناء فيكتوريا يتحول من الأزرق إلى الرمادي إلى الأسود. وفي السادسة صباحًا، بدأت الأبراج المالية على الضفة المقابلة تضيء واحدًا بعد آخر. طرق أحد المساعدين الباب مرة ليسأل عما ينبغي ترتيبه للعشاء. قلت لا داعي.}

\StoryParagraph{3}{لم أبك. وظللت أنتظر أن أبكي.}

\NarrativeBreak

\StoryParagraph{0}{لاحقًا، حين عرض عليّ مساعدي تسوية الحسابات، كان أحد البنود تعويضًا للفندق في تينيريفي. بعد مغادرتنا، ظل الشريط يُلصق عدة أيام أخرى. لم يخبر أحد الفتاتين أن بوسعهما التوقف. قال الفندق إن الشريط ظل مثبتًا مدة أطول مما ينبغي، وإن إلصاقه تكرر أكثر مما ينبغي. ولم يكن من الممكن إزالة بقاياه عن إطار النافذة. فكان لا بد من استبدال لوح الزجاج كله.}

\StoryParagraph{0}{قلت: «موافق».}

\StoryParagraph{3}{كان لوح الزجاج ذاك يواجه المحيط الأطلسي.}

\StoryParagraph{0}{كانت كلفة وقود الطائرة في الفاتورة أعلى من كلفة الطابق بأكمله في مونتاج. سألت عن الأمر. قال مساعدي إن طائرة \textenglish{A330} مصممة لمئتي شخص. ومع وجود ثلاثين منا فقط، كانت أخف مما ينبغي، وكان اتزانها غير مناسب للإقلاع. لذلك كان عليهم في كل مرة أن يضيفوا الوقود حتى تصبح ثقيلة بما يكفي. وحين نصل، لا يمكن إبقاء الوقود الزائد على متنها، بل يجب تفريغه.}

\StoryParagraph{0}{قلت: «في كل مرة؟»}

\StoryParagraph{3}{قال: «في كل مرة».}

\StoryParagraph{0}{وكان هناك أيضًا صندوق أعيد من العيادة في لوس أنجلوس. حقن لتحفيز المبيض تكفي ثمانية أيام. وكان يجب حفظها مبردة. لم تُفتح أي واحدة منها.}

\StoryParagraph{3}{وكان هناك أيضًا عربون وكالة تأجير الأرحام: مئتا ألف دولار، غير قابل للاسترداد. كانت الأم البديلة قد اختيرت في يناير. كانت من كاليفورنيا، في الرابعة والثلاثين، ولديها طفلان. كانت صورتها في الملف. لم أنظر إليها قط. كانت مستعدة منذ يناير.}

\StoryParagraph{3}{بعد أن غادرت جينغ تيان، أخرجت صورة عام 2007 ونظرت إليها مرة. مر تسعة عشر عامًا. ظلت وحدات البكسل في كل جهاز تتحسن، فيما بدا أن وحدات البكسل في تلك الصورة تزداد سوءًا. كان وجهها ضبابيًا. كبّرت الصورة، ونظرت طويلًا، محاولًا العثور على أي فرق بين هيئتها آنذاك وهيئتها لاحقًا. لم أعثر على فرق واحد. فالصورة لم تكن يومًا واضحة بما يكفي لأرى.}

\StoryParagraph{0}{قصصت أظافري ذات يوم. وحين انتهيت، مررت إصبعي عليها. كانت حادة.}

\StoryParagraph{0}{تذكرت أن ذلك الشيء لا يزال في الدرج، الشيء الذي تركته وراءها.}

\StoryParagraph{0}{أخرجته ومررته على أظافري مرتين. لم يبد الأمر صحيحًا. كانت قد قالت إن عليّ تحريكه في اتجاه واحد. وكنت قد نسيت أي اتجاه.}

\StoryParagraph{0}{أعدته إلى مكانه.}

\StoryParagraph{3}{ما زالت أظافري تخدش الناس.}

\StoryParagraph{0}{ظهر تذكير على هاتفي ذات صباح. كان مساعدي قد ضبطه في العام السابق أثناء ترتيب الجدول: قبل موعد الولادة بثلاثة أشهر، جهّز الغرفة.}

\StoryParagraph{3}{مررت إصبعي على الشاشة فأزلته.}

\StoryParagraph{0}{منذ أن غادرت، وأنا أفكر في سؤال واحد.}

\StoryParagraph{0}{لو أنني أعطيتها إياه في ذلك اليوم، هل كانت ستبقى؟}

\StoryParagraph{0}{فكرت في الأمر طويلًا.}

\StoryParagraph{3}{ثم أخرجت الملف لاحقًا وقرأته مرة أخرى. أربعون صفحة ونيف.}

\StoryParagraph{0}{ثم تذكرت الألعاب النارية فوق ميناء فيكتوريا ليلة عيد الربيع، كما بدت من الطابق الأعلى.}

\StoryParagraph{0}{اندفع مليون شخص كمد البحر. انتهت الألعاب النارية. وانحسر مليون شخص.}

\StoryParagraph{0}{مليون شخص.}

\StoryParagraph{0}{كل عام سواء.}

\StoryParagraph{0}{هونغ كونغ في الثالثة صباحًا:}

\StoryParagraph{3}{لم يحدث شيء.}
